\chapter{Experiments}
\label{ch:exp}

This chapter evaluates SEER on four open-domain QA benchmarks against published adaptive-RAG baselines. Section~\ref{sec:exp-setup} describes datasets, metrics, and the alignment protocol with prior work. Section~\ref{sec:main-results} reports the main results. Sections~\ref{sec:ablation-tau}--\ref{sec:ablation-train} present ablations on the gating threshold, inference mode, decomposer training data, and synthesis prompt. Section~\ref{sec:trigger-analysis} analyses per-step retrieval trigger rates, and Section~\ref{sec:breakdown} provides per-type and per-hop breakdowns.

\section{Experimental Setup}
\label{sec:exp-setup}

\subsection{Datasets}

Four QA benchmarks are evaluated, three multi-hop and one single-hop (Table~\ref{tab:datasets} in Chapter~\ref{ch:related}). To align with the DeepRAG~\cite{deeprag} evaluation protocol, we use the \textbf{first 1{,}000 questions of each dev set} (no shuffle), which exactly matches the sample slicing implemented in DeepRAG's public release. For MuSiQue, we additionally report full-dev results ($n = 2{,}417$) as a stricter evaluation.

\subsection{Retrieval Corpus}

We follow the \emph{per-question candidate paragraph} retrieval protocol used by Adaptive-RAG~\cite{adaptiverag}, IRCoT~\cite{ircot}, and the original benchmark releases: each question carries $\sim$20 candidate paragraphs (2 gold + 18 distractors for HotpotQA / 2WikiMultihopQA; MuSiQue ships its own distractor pool; PopQA passages are retrieved from a Wikipedia dump pre-filtered by entity match). BM25 ranks within this per-question pool. DeepRAG instead uses a global Wikipedia BM25 index over a DPR dump; both protocols are reported in their original papers, and the per-question setup is the more common choice in the adaptive-RAG literature.

\subsection{Metrics}

We report SQuAD-style \textbf{exact match (EM)} and token-level \textbf{F1} after standard normalisation (lowercasing, article removal, punctuation stripping). When the dataset provides multiple gold aliases (e.g.\ PopQA, NQ), we take the maximum over aliases. Additionally we report:
\begin{itemize}
    \item \textbf{Retrieval rate (R\%)} --- the fraction of sub-question steps that trigger retrieval, measured as $\sum_q |\{i : \pi(s_i) = 1\}| \,/\, \sum_q |\mathcal{D}(q)|$.
    \item \textbf{Decomposition steps} --- the average number of sub-questions emitted per query.
\end{itemize}

\subsection{Baselines}

We compare against six published methods plus internal control baselines:

\begin{itemize}
    \item \textbf{No-RAG CoT.} Chain-of-thought without retrieval (parametric memory only).
    \item \textbf{Single-RAG.} One BM25 retrieval pass on the original query, then standard reader.
    \item \textbf{FLARE}~\cite{flare}. Sentence-level lookahead probe; reported numbers from \cite{deeprag}.
    \item \textbf{DRAGIN}~\cite{dragin}. Entropy + self-attention gate; reported numbers from \cite{deeprag}.
    \item \textbf{UAR}~\cite{deeprag}. Reported in DeepRAG comparison.
    \item \textbf{Adaptive-RAG}~\cite{adaptiverag}. T5-Large query classifier.
    \item \textbf{DeepRAG}~\cite{deeprag}. SFT + binary tree search + DPO (paper numbers and our 8B re-run).
    \item \textbf{Ours, neg\_max\_ent} (internal control). Pre-SEER token-entropy gate without decomposition.
\end{itemize}

\subsection{Implementation Details}

All experiments use Llama-3-8B-Instruct as the base model with a LoRA-rank-16 decomposer adapter. Unless stated otherwise, the \textbf{main results use the three-source (12K) decomposer} trained on MuSiQue + 2WikiMultihopQA + PopQA (Chapter~\ref{ch:method}), because it yields a \emph{single unified model} that performs acceptably across both single- and multi-hop query distributions without any per-dataset switching --- the central robustness claim of this work. A two-source (8K, MuSiQue + 2WikiMultihopQA) variant, which trades single-hop coverage for stronger multi-hop accuracy, is reported as an ablation in Section~\ref{sec:ablation-train}. Retrieval is BM25-Okapi with $k = 5$. The unified threshold is $\tau = 0.3$ on max-entropy, and all main results use \textbf{Mode C} (LoRA adapter always on). Greedy decoding (\texttt{do\_sample=False}) is used throughout for reproducibility.

\section{Main Results}
\label{sec:main-results}

Table~\ref{tab:main-results} compares SEER against all baselines on the aligned 1{,}000-question dev slice.

\begin{table}[h]
    \centering
    \caption[Main Results: EM / F1 on 1000-question Dev Slice]{Main results (EM / F1 \%) on the first 1{,}000 questions of each dev set. SEER uses the three-source (12K) decomposer with synthesis prompt v2. Best result per column in \textbf{bold}; second-best \underline{underlined}.}
    \label{tab:main-results}
    \setlength{\tabcolsep}{4pt}
    \small
    \begin{tabular}{ l cc cc cc cc }
        \toprule
        \multirow{2}{*}{Method} & \multicolumn{2}{c}{HotpotQA} & \multicolumn{2}{c}{2Wiki} & \multicolumn{2}{c}{MuSiQue} & \multicolumn{2}{c}{PopQA} \\
        \cmidrule(lr){2-3} \cmidrule(lr){4-5} \cmidrule(lr){6-7} \cmidrule(lr){8-9}
        & EM & F1 & EM & F1 & EM & F1 & EM & F1 \\
        \midrule
        \multicolumn{9}{l}{\textit{Reported in DeepRAG paper}~\cite{deeprag}} \\
        CoT (no-RAG)              & 27.2 & 37.8 & 28.2 & 34.9 & 13.7 & 23.0 & --   & --   \\
        FLARE~\cite{flare}        & 23.4 & 32.9 & 30.5 & 37.5 & 17.6 & 23.6 & --   & --   \\
        DRAGIN~\cite{dragin}      & 27.6 & 38.1 & 29.1 & 35.7 & 11.8 & 20.0 & --   & --   \\
        UAR                       & 29.7 & 40.7 & 34.8 & 42.4 & 19.1 & 28.4 & --   & --   \\
        DeepRAG                   & \textbf{40.7} & \textbf{51.5} & \underline{48.1} & \underline{53.3} & 22.5 & \underline{30.4} & \textbf{42.5} & 47.8 \\
        \midrule
        \multicolumn{9}{l}{\textit{Our re-implementations on Llama-3-8B}} \\
        CoT (no-RAG)              & 25.2 & 35.6 & 31.2 & 39.3 & 11.5 & 20.3 & 20.4 & 23.1 \\
        Single-RAG                & 32.4 & 44.0 & 37.9 & 47.0 & 19.2 & 29.0 & 41.6 & \underline{51.3} \\
        neg\_max\_ent (no decomp) & 32.7 & 44.3 & 36.2 & 45.4 & 19.0 & 29.1 & \underline{42.4} & \textbf{51.6} \\
        DeepRAG-8B + SC           & \underline{39.7} & \underline{50.4} & \textbf{51.1} & \textbf{55.9} & \underline{22.8} & 30.2 & --   & --   \\
        \midrule
        \textbf{SEER (Ours, 12K)} & 36.5 & 49.1 & 37.7 & 46.1 & \textbf{46.2} & \textbf{55.1} & 39.8 & 48.3 \\
        \bottomrule
    \end{tabular}
\end{table}

\subsection{Key Findings}

\textbf{Finding 1: MuSiQue is a decisive win.} SEER reaches 46.2 EM / 55.1 F1 on MuSiQue, exceeding the DeepRAG paper by $+23.7$ EM and $+24.7$ F1, and exceeding every other baseline by $\geq +23$ EM. This is the largest single-dataset margin among the methods compared. MuSiQue is the hardest of the four benchmarks (deeper compositional chains, harder distractors), and SEER's KG-triple decomposition appears particularly well-suited to its structured composition.

\textbf{Finding 2: a single unified model is competitive everywhere, including single-hop.} This is the central result. The same 12K decomposer that wins MuSiQue also reaches 39.8 EM / 48.3 F1 on the single-hop PopQA benchmark --- within $2.7$ EM of the specialised DeepRAG ($42.5$) and within $3.3$ F1 of the strongest no-decomposition baseline (neg\_max\_ent, $51.6$) --- \emph{without any per-dataset switching or routing}. No prior method in the comparison spans the single-hop / multi-hop divide with one model: DeepRAG reports PopQA only with its single-hop configuration, while our own multi-hop pipelines collapse to $14.2$ EM under the two-source decomposer (Section~\ref{sec:ablation-train}). SEER's robustness across query distributions, at a fixed inference policy, is the contribution that the headline numbers should be read against.

\textbf{Finding 3: the cost of unification falls on 2Wiki and HotpotQA EM.} Relative to a multi-hop-specialised decomposer (8K, Section~\ref{sec:ablation-train}), the unified model gives up EM on 2Wiki ($45.0 \to 37.7$) and HotpotQA ($39.6 \to 36.5$), and consequently trails DeepRAG on these two. Part of this gap is span-trimming (factually correct spans carrying extraneous modifiers, addressed by synthesis prompt v2 in Section~\ref{sec:synth-prompt}) and part is genuine under-decomposition induced by the single-hop training signal. This trade-off is quantified in Section~\ref{sec:ablation-train} and discussed as the principal open problem in Chapter~\ref{ch:discussion}.

\section{Ablation: Gating Threshold $\tau$}
\label{sec:ablation-tau}

The threshold $\tau$ controls the retrieval / parametric trade-off. Table~\ref{tab:tau-sweep} sweeps $\tau \in [0, \infty)$ on MuSiQue (200-question sample) with the max-entropy signal.

\begin{table}[h]
    \centering
    \caption[Threshold Sweep on MuSiQue 200]{Threshold sweep of $\tau$ on max-entropy signal (MuSiQue, $n=200$). Retrieval\% is the fraction of sub-question steps that trigger retrieval.}
    \label{tab:tau-sweep}
    \setlength{\tabcolsep}{10pt}
    \begin{tabular}{ c c c c l }
        \toprule
        $\tau$ & Retrieval\% & EM & F1 & Notes \\
        \midrule
        0.0 (always)    & 100.0 & 47.5 & 58.9 & Upper bound \\
        0.2             & 94.0  & 47.5 & 58.9 & Gate matches always-RAG \\
        \textbf{0.3}    & 90.8  & 47.0 & 58.8 & \textbf{Sweet spot} \\
        0.5             & 87.5  & 46.5 & 58.0 & $-1$ EM \\
        0.8             & 78.1  & 41.5 & 54.3 & Cliff begins \\
        1.0             & 75.1  & 41.0 & 53.6 & $-6$ EM \\
        1.5             & 65.8  & 36.5 & 51.3 & $-11$ EM \\
        2.5             & 50.1  & 27.0 & 41.6 & $-20$ EM \\
        3.5             & 26.7  & 18.5 & 30.1 & $-29$ EM \\
        $\infty$ (never) & 0.0  & 7.5  & 18.2 & Lower bound \\
        \bottomrule
    \end{tabular}
\end{table}

The curve exhibits a sharp \textbf{cliff above $\tau = 0.5$}: as long as retrieval rate stays $\geq 80\%$, EM is within 1 point of the always-retrieve upper bound, while below this rate accuracy degrades almost linearly. The cliff position is consistent with the pre-SEER \texttt{neg\_max\_ent@-0.80} observation that the entropy distribution of correct-answer probes concentrates near zero. This justifies the choice of $\tau = 0.3$ as a single dataset-agnostic threshold --- it sits well inside the safe region while still saving 7--9 percentage points of retrieval cost (Section~\ref{sec:trigger-analysis}).

\section{Ablation: Inference Mode B vs C}
\label{sec:ablation-mode}

Table~\ref{tab:mode} compares Mode B (LoRA off during answering) and Mode C (LoRA on throughout) on MuSiQue 200.

\begin{table}[h]
    \centering
    \caption[Mode B vs Mode C]{Inference mode ablation (MuSiQue, $n=200$).}
    \label{tab:mode}
    \setlength{\tabcolsep}{10pt}
    \begin{tabular}{ l l c c c }
        \toprule
        Mode & Description & EM & F1 & Retrieval\% \\
        \midrule
        B & Decompose / synth with LoRA; answer with base model & 35.0 & 47.6 & 71.1 \\
        \textbf{C} & LoRA enabled at all stages & \textbf{40.5} & \textbf{53.3} & 69.6 \\
        \bottomrule
    \end{tabular}
\end{table}

Mode C wins by $+5.5$ EM and $+5.7$ F1 with effectively identical retrieval rate. The decomposition adapter does not degrade answer-only generation; on the contrary, it slightly improves it, likely because the adapter has been exposed to entity-property patterns during training that overlap with the answerer's task. \textbf{All subsequent experiments use Mode C.}

\section{Ablation: Decomposer Training Data}
\label{sec:ablation-train}

SEER's default decomposer is trained on the 12K three-source mix (MuSiQue + 2Wiki + PopQA, ``3DS''). To isolate the effect of the single-hop component, we ablate it: a 8K two-source decomposer (MuSiQue + 2Wiki only) is trained under identical hyper-parameters. Table~\ref{tab:train-data-ablation} reports both.

\begin{table}[h]
    \centering
    \caption[Training Data Composition Ablation]{Effect of removing the 4K PopQA single-hop component from the decomposer training corpus. All values are EM / F1 on 1000-question dev slices. The 12K three-source mix is SEER's default.}
    \label{tab:train-data-ablation}
    \setlength{\tabcolsep}{6pt}
    \begin{tabular}{ l c c c c }
        \toprule
        Training mix & HotpotQA & 2Wiki & MuSiQue & PopQA \\
        \midrule
        \textbf{12K three-source (default)} & 36.5 / 49.1 & 37.7 / 46.1 & 46.2 / 55.1 & 39.8 / 48.3 \\
        8K two-source (no PopQA)            & 39.6 / 52.4 & 45.0 / 53.7 & 48.0 / 57.1 & 14.2 / 21.8 \\
        $\Delta$ (12K $-$ 8K)               & $-3.1$ / $-3.3$ & $-7.3$ / $-7.6$ & $-1.8$ / $-2.0$ & $\textbf{+25.6}$ / $\textbf{+26.5}$ \\
        \bottomrule
    \end{tabular}
\end{table}

The single-hop component is what makes SEER a unified model: it lifts PopQA by $+25.6$ EM (from a broken $14.2$ to a competitive $39.8$) by teaching the decomposer that a \emph{one-step plan is a valid output}, directly mitigating the over-decomposition failure mode. The cost is multi-hop accuracy: 2Wiki loses $-7.3$ EM, HotpotQA $-3.1$, and MuSiQue $-1.8$, because the same single-hop signal also pushes the decomposer to under-decompose genuine 2-hop questions. MuSiQue is the most robust to the change --- its chains are deeper and the KG-triple supervision is internally consistent. This is a clear example of negative transfer: \textbf{single-hop and multi-hop decomposition supervision interfere.} SEER adopts the 12K mix as default because the robustness of a single model across the full query distribution outweighs the multi-hop EM it costs; closing that gap --- e.g.\ by conditioning the decomposer on a length-target token, analogous to length control in summarisation, so one model can emit different decomposition depths on demand --- is the principal future direction (Chapter~\ref{ch:discussion}).

\section{Ablation: Synthesis Prompt}
\label{sec:synth-prompt}

The terse-answer rules and worked few-shot examples in the synthesis prompt (Section~\ref{sec:synthesize}) were added in response to the observed EM/F1 gap on 2Wiki bridge-comparison questions. The change is \emph{prompt-only} --- no re-training --- so its effect is independent of the decomposer mix; the measurements below were taken on the two-source decomposer during development and transfer to the default model. PopQA is omitted: it has a single sub-step, so there is nothing to synthesise and the prompt has no effect.

\begin{table}[h]
    \centering
    \caption[Synthesis Prompt Ablation]{Effect of upgrading the synthesis prompt from v1 (rules only) to v2 (rules + comparison + Yes/No few-shot), measured on the multi-hop benchmarks. EM \% on 1000-question slices.}
    \label{tab:synth-prompt}
    \setlength{\tabcolsep}{12pt}
    \begin{tabular}{ l c c c }
        \toprule
        Prompt & HotpotQA & 2Wiki & MuSiQue \\
        \midrule
        v1 (rules only)            & 37.8 & 35.9 & 45.9 \\
        v2 (rules + few-shot)      & 39.6 & 45.0 & 48.0 \\
        $\Delta$                   & $+1.8$ & $\textbf{+9.1}$ & $+2.1$ \\
        \bottomrule
    \end{tabular}
\end{table}

The 2Wiki gain ($+9.1$ EM) is the headline: it derives almost entirely from the bridge-comparison subtype, where the original synthesizer emitted KG-triple-shaped strings (e.g.\ \texttt{X >> date of birth}) instead of resolving the comparison. The few-shot examples teach the model to map ``compare two dates'' to ``return the entity owning the later date.'' MuSiQue and HotpotQA also benefit moderately. \textbf{This is a free win: no re-training, no regression, every multi-hop benchmark improved.}

\section{Retrieval Trigger Rate Analysis}
\label{sec:trigger-analysis}

Table~\ref{tab:trigger} reports per-dataset retrieval trigger rates under the default $\tau = 0.3$.

\begin{table}[h]
    \centering
    \caption[Per-Dataset Retrieval Trigger Rate]{Retrieval rate at $\tau = 0.3$ (max-entropy gate). $|\mathcal{D}|$ is average decomposition length.}
    \label{tab:trigger}
    \setlength{\tabcolsep}{8pt}
    \begin{tabular}{ l c c c }
        \toprule
        Dataset & $|\mathcal{D}|$ (avg steps) & Retrieve / Total steps & Trigger Rate \\
        \midrule
        HotpotQA  & 2.26 & 2060 / 2257 & 91.3\% \\
        2Wiki     & 2.47 & 2357 / 2472 & 95.3\% \\
        MuSiQue   & 2.03 & 1832 / 2032 & 90.2\% \\
        PopQA     & 1.00 &  993 / 1000 & 99.3\% \\
        \bottomrule
    \end{tabular}
\end{table}

Three observations:
\begin{enumerate}
    \item \textbf{Trigger rates cluster between 89--95\% for multi-hop benchmarks.} The gate does not skip aggressively, which is consistent with the threshold-sweep cliff at 80\% (Section~\ref{sec:ablation-tau}): SEER deliberately operates in the safe region above the cliff.
    \item \textbf{2Wiki has the highest trigger rate (95.3\%)}, reflecting that its Wikidata-derived questions reference long-tail entities the base Llama-3 has weak parametric coverage of.
    \item \textbf{PopQA triggers retrieval on nearly every step (99.3\%)} because its questions hinge on long-tail entity facts where parametric memory is unreliable. The high trigger rate is not the cause of PopQA's low EM --- the cause is over-decomposition before retrieval is even invoked.
\end{enumerate}

The gate therefore saves 5--11 percentage points of retrieval calls compared to always-retrieve, while losing $\leq 0.5$ EM (Table~\ref{tab:tau-sweep}) --- a clear Pareto improvement.

\section{Per-Type and Per-Hop Breakdowns}
\label{sec:breakdown}

\subsection{2Wiki by Question Type}

Table~\ref{tab:2wiki-type} breaks down 2Wiki performance by official question type, exposing where the bulk of the synthesis-prompt gain came from.

\begin{table}[h]
    \centering
    \caption[2Wiki Per-Type Breakdown]{2Wiki per-type EM (1000-question slice, $\tau=0.3$, synthesis prompt v1 vs v2).}
    \label{tab:2wiki-type}
    \setlength{\tabcolsep}{10pt}
    \begin{tabular}{ l r c c c }
        \toprule
        Type & $n$ & v1 EM & v2 EM & $\Delta$ \\
        \midrule
        inference          & 128 & 58.6 & 60.2 & $+1.6$ \\
        compositional      & 428 & 46.3 & 50.0 & $+3.7$ \\
        comparison         & 216 & 29.6 & 41.7 & $+12.1$ \\
        bridge\_comparison & 228 & ~9.6 & 39.5 & $\textbf{+29.9}$ \\
        \midrule
        Total              & 1000 & 35.9 & 45.0 & $+9.1$ \\
        \bottomrule
    \end{tabular}
\end{table}

Bridge-comparison (4-step double-chain questions) went from \textbf{9.6\% to 39.5\% EM}, a near 4$\times$ improvement attributable entirely to the synthesis prompt v2. Inference and compositional questions, which were already working, see modest gains.

\subsection{MuSiQue by Hop Depth}

MuSiQue's full dev set (2{,}417 questions) is partitioned by hop depth in Table~\ref{tab:musique-hop}.

\begin{table}[h]
    \centering
    \caption[MuSiQue Per-Hop EM]{MuSiQue full-dev EM by hop class (2{,}417 questions, $\tau=0.3$).}
    \label{tab:musique-hop}
    \setlength{\tabcolsep}{10pt}
    \begin{tabular}{ l r c }
        \toprule
        Hop class & $n$ & EM \\
        \midrule
        2hop   & 1252 & \textbf{43.0} \\
        3hop\_v1 & 568  & 36.6 \\
        3hop\_v2 & 192  & 15.1 \\
        4hop\_v1 & 246  & 25.6 \\
        4hop\_v2 & 64   & 14.1 \\
        4hop\_v3 & 95   & ~9.5 \\
        \midrule
        Total   & 2417 & 35.4 \\
        \bottomrule
    \end{tabular}
\end{table}

Performance is monotonically decreasing in hop count, as expected. The 4hop variants are the principal short-coming and align with the decomposer count-match ablation in Chapter~\ref{ch:method}: the decomposer's count-match drops to 40--46\% on 4-hop chains, propagating errors into both retrieval and synthesis. The first 1{,}000 questions of MuSiQue dev consist entirely of 2-hop examples (the dataset ships sorted by hop count), which is why the aligned-protocol number (48.0 EM) is higher than the full-dev number (35.4 EM); we report both in the interest of fairness.

\section{Summary}

The experiments support three conclusions:
\begin{itemize}
    \item \textbf{One unified model spans single- and multi-hop.} The default three-source (12K) SEER decisively wins MuSiQue ($+23.7$ EM over DeepRAG) while remaining competitive on single-hop PopQA ($39.8$ EM, within $2.7$ of specialised DeepRAG; $48.3$ F1, above it) --- with no per-dataset routing. This robustness across the query distribution is the central contribution.
    \item \textbf{The gating signal is robust.} A single threshold $\tau = 0.3$ on max-entropy works across all four datasets without per-dataset tuning, and the threshold-sweep cliff at 80\% retrieval rate justifies the choice.
    \item \textbf{Unification has a measurable cost on 2Wiki / HotpotQA EM.} Against a multi-hop-only (8K) decomposer, adding single-hop supervision rescues PopQA by $+25.6$ EM but costs $7.3$ EM on 2Wiki and $3.1$ on HotpotQA --- a clear sign of supervision interference. Closing this gap with a single model is the principal open problem, discussed at length in Chapter~\ref{ch:discussion}.
\end{itemize}
